\section{Conclusions \& Discussion}
\label{sec:conclusion}

\paragraph{The take.} Hyperbolic geometry helps with one specific kind
of tree-aware behaviour --- preserving local neighborhoods, so that
the nearest neighbors of an artwork share its sub-style --- and it
does not give us a free upgrade on classification or on global tree
fidelity. The right framing is that geometry is a knob with a clear
trade-off, not a hyperparameter to tune to a single optimum.
Practitioners who care about retrieval inside fine-grained taxonomies
should pick the Poincar\'e ball; practitioners who care about getting
the label right on a single image should stay Euclidean. The two
goals are in genuine tension.

\paragraph{What the numbers do not say.} They do not say hyperbolic
embeddings are universally preferable for hierarchical image data ---
our results say the opposite on global metrics. They do not say the
hand-built tree is the \emph{correct} style hierarchy; tree-Spearman
and dendrogram F1 are anchored to that specific taxonomy and a
different art-historian would draw it differently. They do not
generalise beyond medium-scale, Western-canon-heavy WikiArt with a
frozen CLIP backbone --- a different dataset or a fine-tuned encoder
could change the picture meaningfully.

\paragraph{Limitations and the natural extensions.} Three are most
honest to call out. First, only the prototypes live on the manifold;
the head MLP is still Euclidean, so the hyperbolic geometry can only
shape the decision boundary, not the representation itself. A
genuinely hyperbolic head along the lines of M\"obius layers
\citep{ganea2018hyperbolic} would let the geometry do more work and
might shift the trade-off. Second, the regularizer is a single fixed
loss formulation: scale-invariant MSE on pairwise distances.
Alternatives that match the rank order of tree distances (a Spearman
surrogate) or that weight nearby pairs more heavily (Sammon-style)
might pay smaller classification costs to achieve the same tree
fidelity. Third, the ground-truth tree is fixed by hand; a learned
tree along the lines of \citet{chami2020trees} would test whether the
hand-built taxonomy is itself a bottleneck for the global metrics, and
in particular whether the local-vs-global gap is partly an artifact of
asking the embedding to match a tree that the data itself does not
quite support.
